% LockAI 中文论文 / 报告模板（xelatex）
% 用法：复制到 work/ 改正文，然后 `latexmk -xelatex -interaction=nonstopmode paper.tex`
% 要同时交 Word：`todocx paper.tex -o /home/user/outputs/题目.docx`（正文只用下面出现过的写法，Word 版才转得干净）
\documentclass[UTF8,zihao=-4,a4paper,fontset=none]{ctexart}

% ---------- 字体：思源宋体 / 思源黑体 + Times 风格西文 ----------
\setCJKmainfont{Noto Serif CJK SC}[BoldFont=Noto Serif CJK SC Bold]
\setCJKsansfont{Noto Sans CJK SC}[BoldFont=Noto Sans CJK SC Bold]
\setCJKmonofont{Noto Sans Mono CJK SC}
\setCJKfamilyfont{zhsong}{Noto Serif CJK SC}
\setCJKfamilyfont{zhhei}{Noto Sans CJK SC}
\setmainfont{TeX Gyre Termes}
\setsansfont{TeX Gyre Heros}

% ---------- 版面：A4，页边距同 Word 默认，1.5 倍行距 ----------
\usepackage[top=2.54cm,bottom=2.54cm,left=3.17cm,right=3.17cm]{geometry}
\linespread{1.6}
\usepackage{amsmath,amssymb}
\usepackage{graphicx}
\usepackage{booktabs}
\usepackage{caption}
\captionsetup{font=small,labelsep=quad}
\usepackage[hidelinks]{hyperref}

% ---------- 标题样式（黑体，自动编号 1 / 1.1 / 1.1.1，Word 版编号一致） ----------
\ctexset{
  section = {format = \sffamily\zihao{-3}\bfseries, beforeskip = 1.2ex plus .2ex, afterskip = 1ex},
  subsection = {format = \sffamily\zihao{4}\bfseries, beforeskip = 1ex, afterskip = .6ex},
  subsubsection = {format = \sffamily\zihao{-4}\bfseries, beforeskip = .8ex, afterskip = .4ex},
  abstractname = {摘\quad 要},
}
\renewenvironment{abstract}{%
  \par\begin{center}\sffamily\zihao{4}\bfseries\abstractname\end{center}\par\zihao{-4}%
}{\par\medskip}

% ---------- 页码居中 ----------
\pagestyle{plain}

\title{\sffamily\zihao{2}\bfseries 这里写论文题目}
\author{\zihao{-4}作者姓名\quad 学号 2025xxxxxx\\ \zihao{5}某某大学某某学院}
\date{\zihao{5}2026 年 9 月}

\begin{document}
\maketitle

\begin{abstract}
这里写摘要，200~300 字，概括研究问题、方法、主要观点和结论。

\noindent\textbf{关键词：}关键词一；关键词二；关键词三
\end{abstract}

\section{引言}
正文段落直接写，段首自动缩进两字。引用文献用 \verb|\cite{key}|，如已有研究指出这一点\cite{marx1972}。

\section{第一部分}
\subsection{小节标题}
行内公式 $E = mc^2$，行间公式：
\begin{equation}
  f(x) = \int_{0}^{x} e^{-t^2}\,\mathrm{d}t
\end{equation}

表格用三线表：
\begin{table}[htbp]
  \centering
  \caption{表格标题}
  \begin{tabular}{lcc}
    \toprule
    项目 & 数值 A & 数值 B \\
    \midrule
    第一行 & 1.0 & 2.0 \\
    第二行 & 3.0 & 4.0 \\
    \bottomrule
  \end{tabular}
\end{table}

图片：
% \begin{figure}[htbp]
%   \centering
%   \includegraphics[width=.7\linewidth]{figure.png}
%   \caption{图片标题}
% \end{figure}

\section{结论}
总结全文。

% 参考文献按 GB/T 7714—2015 手写格式，只列查证过的真实文献
\begin{thebibliography}{99}
  \bibitem{marx1972} 恩格斯. 自然辩证法[M]. 北京: 人民出版社, 1971.
\end{thebibliography}

\end{document}
